\chapter{Gaussian, matchgate, and fermionic catalytic models}\label{ch:gaussian}
\index{Gaussian circuit}
\index{matchgate}
\index{Bogoliubov transformation}
\index{fermionic Gaussian circuit}
\index{symplectic embedding}


\section{Two Gaussian settings}
Quantum Gaussianity occurs \cite{WeedbrookEtAl2012} in two closely related but algebraically different settings.

For $n$ bosonic modes, a Gaussian unitary acts on quadratures
\[
\mathbf r=(q_1,\ldots,q_n,p_1,\ldots,p_n)^T
\]
by an affine symplectic transformation
\[
\mathbf r\longmapsto S\mathbf r+d,
\qquad
S\in\Sp(2n,\R).
\]
For $n$ fermionic modes, a fermionic Gaussian unitary acts on $2n$ Majorana operators by
\[
\gamma\longmapsto O\gamma,
\qquad O\in SO(2n).
\]
Nearest-neighbor matchgates give a circuit presentation of a fermionic Gaussian subgroup on qubits.

In both cases one can choose a restricted subgroup preserving a flag, mode ordering, particle number, or bipartition, and compare it with a larger Gaussian or non-Gaussian gate set.

\section{An inherited symplectic shear theorem}
Every invertible linear transformation $A\in\GL(n,R)$ has a canonical symplectic embedding
\begin{equation}\label{eq:symplectic-embedding}
\iota(A)=
\begin{pmatrix}
A&0\\
0&A^{-T}
\end{pmatrix}.
\end{equation}
Indeed, with
\[
J=
\begin{pmatrix}0&I\\-I&0\end{pmatrix},
\]
one checks
\[
\iota(A)^T J\iota(A)=J.
\]
Thus every row-transvection identity in the established linear theory induces an exact identity between symplectic shear circuits.

\begin{provedbox}
Under a Gaussian gate alphabet in which each embedded elementary transvection $\iota(I+aE_{ij})$ is one primitive and direction is the only restriction, the universal six-row catalytic atom embeds into a twelve-quadrature Gaussian atom with restricted length seven and ambient length six.
\end{provedbox}

This is a formal inheritance statement for the declared shear alphabet.  It does not assert optimality after arbitrary beam splitters, squeezers, Fourier transforms, or local Gaussian gates are admitted at unit cost.

\begin{figure}[ht]
\centering
\includegraphics[width=.94\textwidth]{gaussian_shear_catalysis.pdf}
\caption{The row-transvection atom embedded as a Gaussian symplectic shear circuit.}
\label{fig:gaussian-shear}
\end{figure}

\section{Bosonic gate hierarchies}
Natural nested models include
\[
\text{triangular shears}
\subset
\text{all Gaussian unitaries}
\subset
\text{Gaussian plus a non-Gaussian primitive}
\subset
U(\mathcal H).
\]
Typical primitive resources are:
\begin{itemize}
\item SUM or controlled displacement gates,
\item beam splitters,
\item single-mode and two-mode squeezers,
\item Fourier transforms,
\item cubic phase gates,
\item photon counting or non-Gaussian state injection.
\end{itemize}
Each inclusion supports a different defect.  A target Gaussian transformation might be shortened by leaving a triangular decomposition but remaining Gaussian.  A more radical question asks whether temporary non-Gaussianity can shorten a final Gaussian target.

\section{Bloch--Messiah versus catalytic synthesis}
A bosonic Gaussian unitary admits a Bloch--Messiah decomposition
\[
S=K_1 D K_2,
\]
where $K_1,K_2$ are passive symplectic transformations and $D$ is a diagonal squeezing layer.  This normal form optimizes conceptual structure, not necessarily a weighted gate count on a specific architecture.

Catalytic synthesis asks whether a short excursion into a basis where many shear or squeezing obligations align can reduce total cost.  The opening and closing passive layers play the role of the catalyst.

\begin{skeletonbox}
A Gaussian catalytic theorem requires:
\begin{enumerate}
\item a declared primitive library and cost vector;
\item a target subgroup, such as triangular, number conserving, or local;
\item an explicit sandwich $KHK^{-1}$ returning to the target subgroup;
\item a lower bound based on symplectic rank, singular values, graph cuts, or mode support;
\item treatment of approximation error when continuous parameters are synthesized numerically.
\end{enumerate}
\end{skeletonbox}

\section{Fermionic Gaussian and matchgate circuits}
Let $c_j,c_j^\dagger$ be fermionic creation and annihilation operators.  Quadratic Hamiltonians generate Gaussian unitaries, and their action is determined by an orthogonal transformation on Majoranas.  Number-conserving Gaussian operations correspond to a unitary action on modes, while general Gaussian operations allow particle-hole mixing.

This yields the hierarchy
\[
\begin{gathered}
U(n)\text{ mode transformations}
\subset SO(2n)\text{ fermionic Gaussian transformations}\\
\subset \text{parity-preserving unitaries}.
\end{gathered}
\]
A temporary Bogoliubov transformation can violate number conservation while preserving parity.  The target may nevertheless return to the number-conserving subgroup.

\begin{definition}[Bogoliubov catalytic defect]
For a number-conserving Gaussian target $U$, compare the minimum cost using number-conserving primitives with the minimum cost using general Gaussian primitives, requiring the endpoint to conserve number exactly.
\end{definition}

\section{Matchgates and geometry}
Nearest-neighbor matchgate circuits \cite{Valiant2002,JozsaMiyake2008} have strong classical simulation properties.  Their geometry changes when SWAP, fermionic SWAP, long-range connectivity, or a single non-matchgate resource is added.  Candidate restricted/ambient pairs include:
\begin{enumerate}
\item nearest-neighbor matchgates versus all-to-all matchgates;
\item number-conserving matchgates versus general matchgates;
\item matchgates versus parity-preserving two-qubit gates;
\item Gaussian circuits versus Gaussian circuits with one reusable non-Gaussian ancilla.
\end{enumerate}

A compiler application is immediate: a temporary mode permutation or particle-hole rotation may allow a batch of otherwise expensive interactions to share one routing or basis-change layer.

\section{Covariance-matrix lower bounds}
For a Gaussian state, the covariance matrix $\Gamma$ transforms as
\[
\Gamma\longmapsto S\Gamma S^T
\]
in the bosonic case and
\[
\Gamma\longmapsto O\Gamma O^T
\]
in the fermionic case.  Possible restricted lower bounds include:
\begin{itemize}
\item rank change of off-diagonal covariance blocks,
\item symplectic singular-value obligations,
\item growth of mode support,
\item graph bandwidth and cut rank,
\item number of nontrivial canonical angles between invariant subspaces.
\end{itemize}
One elementary gate changes only a controlled portion of these data.  A useful invariant must also charge preparation of a carrier that makes progress on many cuts simultaneously.

\section{Temporary non-Gaussianity}
Suppose a non-Gaussian primitive $g$ is introduced and later removed:
\[
U=g h_m\cdots h_1 g^{-1},
\qquad U\in G_{\mathrm{Gauss}}.
\]
If $g h_i g^{-1}$ generates several useful Gaussian effects, the excursion may be catalytic.  However, many standard non-Gaussianity monotones vanish only on states rather than unitaries, and cancellation at the endpoint complicates direct lower bounds.

Candidate monotones and certificates include:
\begin{itemize}
\item degree in the Weyl characteristic function,
\item Wigner negativity generated during the path,
\item operator Schmidt rank across mode cuts,
\item polynomial degree of Heisenberg-evolved quadratures,
\item stabilizer rank for encoded finite-dimensional subspaces.
\end{itemize}

\section{Exact and numerical search}
Finite exact enumeration is available for binary symplectic or orthogonal quotients, but continuous Gaussian groups require a different strategy.
\begin{enumerate}
\item Fix a symbolic gate skeleton and solve polynomial equations for parameters.
\item Use interval arithmetic to certify nonexistence below a length.
\item Exploit canonical decompositions to quotient local or passive layers.
\item For discrete-angle Gaussian libraries, use exact algebraic-number arithmetic.
\item Store both an upper-bound circuit and a dual or algebraic certificate.
\end{enumerate}

\section{Architecture and routing}
On a line or grid, Gaussian and matchgate circuits may be dominated by routing.  Let $G_{\mathrm{conn}}$ be the hardware graph.  A mode permutation $\pi$ can be treated as a catalyst if it is opened, many interactions are executed in the reordered frame, and $\pi^{-1}$ restores the logical labeling.  The benefit is
\[
\text{routing saved across consumers}
-
2\,\text{cost}(\pi).
\]
This is the same amortization law as the linear snapshot family, now applied to mode geometry.

\begin{conjecturebox}
There exist number-conserving or nearest-neighbor Gaussian targets whose minimum weighted cost is reduced by a restored Bogoliubov, mode-permutation, or non-Gaussian excursion.  The first exact threshold depends on the primitive library and is not determined here.
\end{conjecturebox}

\section{Research deliverables}
A complete Gaussian extension should produce:
\begin{enumerate}
\item a finite or semialgebraic model with exact endpoint equality;
\item at least one nontrivial lower-bound invariant;
\item a minimal catalytic atom;
\item a weighted hardware interpretation;
\item a tensor or repeated-mode amplification theorem;
\item code using exact arithmetic or certified intervals.
\end{enumerate}
